\documentclass[11pt]{article}

\usepackage[T1]{fontenc}
\usepackage[utf8]{inputenc}
\usepackage{lmodern}
\usepackage[margin=1.1in]{geometry}
\usepackage{amsmath,amssymb,amsthm,mathtools}
\usepackage{enumitem}
\usepackage{microtype}
\usepackage{xcolor}
\usepackage{booktabs}
\usepackage{array}
\usepackage{titlesec}
\usepackage{hyperref}

\titleformat{\section}
  {\normalfont\LARGE\bfseries}{\thesection}{1em}{}
\titleformat{\subsection}
  {\normalfont\Large\bfseries}{\thesubsection}{1em}{}

\hypersetup{
  pdftitle={Collision Ideals and Off-Diagonal Sheets},
  pdfauthor={Chloe van der Vlugt},
  colorlinks=true,
  linkcolor=blue!55!black,
  citecolor=green!40!black,
  urlcolor=blue!60!black
}

\newtheorem{theorem}{Theorem}[section]
\newtheorem{proposition}[theorem]{Proposition}
\newtheorem{lemma}[theorem]{Lemma}
\newtheorem{corollary}[theorem]{Corollary}
\theoremstyle{remark}
\newtheorem*{remark}{Remark}
\theoremstyle{definition}
\newtheorem{definition}[theorem]{Definition}

\DeclareMathOperator{\Spec}{Spec}
\DeclareMathOperator{\Frac}{Frac}
\DeclareMathOperator{\Gal}{Gal}
\DeclareMathOperator{\Hom}{Hom}
\DeclareMathOperator{\Aut}{Aut}
\DeclareMathOperator{\Ram}{Ram}
\DeclareMathOperator{\Norm}{Norm}
\DeclareMathOperator{\Ann}{Ann}
\DeclareMathOperator{\core}{core}

\newcommand{\A}{\mathbb A}
\newcommand{\C}{\mathbb C}
\newcommand{\Obs}{\operatorname{Obs}}
\newcommand{\id}{\mathrm{id}}
\newcommand{\qedbox}{\hfill$\square$}
\newcommand{\leanpath}[1]{\texttt{\nolinkurl{#1}}}

\newenvironment{leanfiles}
  {\par\medskip\begingroup\footnotesize\raggedright\noindent
   \textit{Lean correspondence.}\ }
  {\par\endgroup\medskip}

\title{Collision Ideals and Off-Diagonal Sheets}
\author{Chloe van der Vlugt}
\date{\today}

\begin{document}

\maketitle

\begin{abstract}
We study the scheme-theoretic collision relation of a polynomial self-map,
namely its fiber product over the target, and describe the geometry away
from the diagonal through boundary-supported inertia and Galois-closure
sheets.
For a polynomial map
\(F=(F_1,\dots,F_n)\colon\A^n_{\C}\to\A^n_{\C}\), define in
\(S=\C[x_1,\dots,x_n,y_1,\dots,y_n]\) the collision and diagonal ideals
\[
  I_R=\bigl(F_i(x)-F_i(y)\bigr)_{i=1}^n,
  \qquad
  I_\Delta=(x_i-y_i)_{i=1}^n.
\]
The image of \(I_\Delta\) in the collision algebra \(C_F:=S/I_R\) is the
obstruction ideal
\(\Obs(F):=I_\Delta/I_R\); it records the failure of the collision scheme
\(\Spec C_F\) to coincide with the diagonal.  Its vanishing characterizes
polynomial automorphisms:
\[
  \Obs(F)=0
  \quad\Longleftrightarrow\quad
  F\text{ is a polynomial automorphism}.
\]
In dimension two we give two module-theoretic formulations of the same
fixed-map conclusion.  Planar boundary coherence asks that the boundary
local cohomology
\(H^1_{\partial X}(\overline X,\mathcal O_{\overline X})\) be a finite
\(\overline A\)-module; it forces \(\partial X=\varnothing\), hence makes
the Keller map finite \'etale and therefore an automorphism.  Planar
ramification rigidity is the divisorial criterion that
\(\Omega_{A_N/B}\) have no height-one support, equivalently that it have
finite length.  On the common Galois normalization \(Z=\Spec A_N\), we
construct the pulled-back
affine opens \(U_g\), their canonical boundary ideals
\(\mathfrak j_g=I_{A_N}(Z\setminus U_g)\), and the fixed-locus ideals
\(\mathfrak a_C\).  The height-one dictionary identifies moving-sheet
coverage and finite-group boundary separation as exact criteria for
ramification rigidity, after which purity and finite-\'etale rigidity give
\(N=L=K\) and collision vanishing.  Although \(\overline A\) and \(A_N\)
are already finite over \(B\), neither of these additional finiteness
criteria follows formally.  Universal validity of either criterion is
equivalent to the two-dimensional Jacobian conjecture and is not proved
here.  The boundary divisor-class argument does, however, prove the
normal/Galois case and therefore excludes generic degree two.
For the separable nonnormal cubic class in dimension three, let
\[
  K=\C(F_1,F_2,F_3),
  \qquad
  L=\C(x_1,x_2,x_3),
\]
and let \(N\) be the normal closure of \(L/K\).  Then
\[
  L\otimes_K L\simeq L\times N,
  \qquad
  \Gal(N/K)\simeq S_3.
\]
The generic off-diagonal factor \(N\) is the ordered-root Galois-closure sheet.
Consequently, \(\Obs(F)\neq0\).  The algebraic construction, the
pulled-back boundary dictionary, and the implications from supplied
coverage or separation hypotheses are formalized in Lean~4.  The
boundary-coherence and finite-length support arguments remain explicit
Lean interfaces whose manuscript proofs have not yet been formalized, as
do the standard external geometric statements.
\end{abstract}

\begin{center}
  \small Licensed under
  \href{https://creativecommons.org/licenses/by/4.0/}
    {CC BY 4.0}.
\end{center}

\tableofcontents

\section{Introduction}
\label{sec:introduction}

The Jacobian conjecture asks whether a polynomial self-map of affine space
with constant nonzero Jacobian determinant must be a polynomial
automorphism.  Dimension two remains open.  Rather than treating
injectivity as an external property of the map, this paper studies the
scheme-theoretic fiber product of the map with itself and asks when its
off-diagonal part disappears.

There is one divisorial endgame, expressed in several exact languages.  The
paper does not prove the Jacobian conjecture.  It organizes the remaining
problem around the following two module-theoretic formulations:
\[
  \begin{aligned}
  \operatorname{PBC}(F)
    &:\Longleftrightarrow
      H^1_{\partial X}(\overline X,\mathcal O_{\overline X})
      \text{ is a finite }\overline A\text{-module},\\
  \operatorname{PRR}(F)
    &:\Longleftrightarrow
      \Omega_{A_N/B}\text{ has no height-one support}\notag\\[-2pt]
    &\hspace{4.7em}\Longleftrightarrow
      \Omega_{A_N/B}\text{ has finite length over }A_N.
  \end{aligned}
\]
Planar boundary coherence \(\operatorname{PBC}(F)\) is the corresponding
global properness formulation:
\[
  \operatorname{PBC}(F)
  \Longrightarrow \partial X=\varnothing
  \Longrightarrow f_F\text{ finite \'etale}
  \Longrightarrow F\text{ an automorphism}.
\]
Planar ramification rigidity \(\operatorname{PRR}(F)\) is the
differential-module formulation of the codimension-one premise:
\[
  \operatorname{PRR}(F)
  \Longrightarrow Z/Y\text{ finite \'etale}
  \Longrightarrow N=K
  \Longrightarrow L=K
  \Longrightarrow q_F=0
  \Longrightarrow I_R=I_\Delta.
\]

The divisorial endgame is controlled by an exact codimension-one obstruction on
the Galois normalization.  A ramified divisor \(E\) is eliminated as soon
as one conjugate sheet moved by \(I_E\) still contains its generic point.
In coordinates this is the conjugate-integrality condition
\[
  I_E\neq1
  \Longrightarrow
  \exists g\in G,\quad
  I_E\nsubseteq gHg^{-1},\qquad
  w_E(g(x_1)),w_E(g(x_2))\geq0.
\]
We call this planar moving-sheet coverage.  Its ideal-theoretic form says
that the fixed locus of a prime-order subgroup and the boundaries of all
the sheets it moves have no common height-one component.  Establishing
this universally---equivalently, establishing planar ramification rigidity
universally---is equivalent to \(JC(2)\), so the reformulation locates the
missing theorem without assuming it.

Only after all height-one points have been shown unramified do purity and
finite-\'etale rigidity enter: they make the normal closure trivial and
collapse the collision relation to the diagonal.  This order is essential
to the argument.  The paper also develops boundary divisor-class and
principal-parts criteria that may imply coverage under additional
hypotheses; these are prospective mechanisms, not additional endgames or
proofs of universal coverage.
The divisor-class mechanism does prove unconditionally that a normal---hence
Galois---marked extension \(L/K\) is trivial.  In particular, a hypothetical
planar counterexample must have \(L/K\) nonnormal and generic degree at least
three.

The normalization rings \(\overline A\) and \(A_N\) are finite
\(B\)-modules from the outset.  This makes \(\overline X\) and \(Z\)
affine varieties: they are integral affine schemes of finite type over
\(\C\).  It does not make \(A\) finite over \(B\), make boundary local
cohomology finite, or make \(\Omega_{A_N/B}\) finite length.  Likewise,
the vanishing of ordinary higher quasi-coherent cohomology on an affine
scheme does not imply the vanishing or finiteness of local cohomology with
support in its deleted boundary.

In the separable nonnormal cubic class in dimension three, the opposite
phenomenon is visible generically: the off-diagonal factor is the
ordered-root \(S_3\)-normal-closure sheet.  This gives a structural
comparison with the planar criteria while keeping its assumptions explicit.

Section~\ref{sec:general-construction} develops collision ideals and the
dimension-independent normalization diagram.
Section~\ref{sec:planar-vanishing} constructs the planar pulled-back
boundaries, proves the valuation dictionary, and states the equivalent
coverage and separation criteria before applying the endgame.
Section~\ref{sec:cubic-s3-obstruction} treats the cubic \(S_3\) class.

\section{General construction}
\label{sec:general-construction}

The organizing question is how the failure of a polynomial self-map to be
one-to-one is encoded by the part of its fiber square lying off the diagonal.
The collision ideal first isolates that part scheme-theoretically; the
normalization diagram then describes it through conjugate sheets and inertia.
These constructions are dimension-independent.  Section
\ref{sec:planar-vanishing} specializes them to planar Keller maps and locates
the obstruction at the deleted normalization boundary, while Section
\ref{sec:cubic-s3-obstruction} identifies the generic off-diagonal sheet in
the separable nonnormal cubic class with an \(S_3\)-normal closure.  The two
specializations respectively describe when the collision scheme collapses to
the diagonal and what its surviving generic sheet looks like when it does not.

\subsection{The two ideals and their obstruction}
\label{sec:collision-ideals}

Throughout the paper, the base field is \(\C\).  Let
\[
  F=(F_1,\dots,F_n)\colon
  X=\A^n_{\C}\longrightarrow\A^n_{\C}
\]
be a polynomial map, and let
\[
  S=\C[x_1,\dots,x_n,y_1,\dots,y_n].
\]
The map \(F\) is Keller if
\[
  \det JF\in\C^\times.
\]
Define
\[
  I_R(F)=\bigl(F_i(x)-F_i(y)\bigr)_{i=1}^{n},
  \qquad
  I_\Delta=(x_i-y_i)_{i=1}^{n},
  \qquad
  \Obs(F):=I_\Delta/I_R(F).
\]
The quotient \(\Obs(F)\) is naturally an ideal of the collision ring
\(C_F:=S/I_R(F)\); we call it the obstruction ideal.
For a fixed map \(F\), we abbreviate \(I_R(F)\) to \(I_R\).

\begin{proposition}[Canonical containment]
\label{prop:collision-contained-in-diagonal}
For every polynomial map \(F\),
\[
  I_R(F)\subseteq I_\Delta.
\]
\end{proposition}

\begin{proof}
For \(H\in\C[x_1,\dots,x_n]\), the difference \(H(x)-H(y)\) belongs to
\((x_1-y_1,\dots,x_n-y_n)\).  This follows by induction on \(H\), or by
telescoping each monomial one coordinate at a time.  Apply the statement
to each coordinate \(F_i\).
\end{proof}

Let \(A:=\C[x_1,\dots,x_n]\).  Diagonal evaluation induces a surjective ring
map
\[
  \bar\mu_F\colon C_F\longrightarrow A,
  \qquad
  [s]_{I_R}\longmapsto s(x,x).
\]

\begin{proposition}[The diagonal kernel]
\label{prop:obstruction-kernel}
As ideals of \(C_F\),
\[
  \ker(\bar\mu_F)
  =
  \Obs(F)
  =
  I_\Delta/I_R.
\]
Consequently,
\[
  \ker(\bar\mu_F)=0
  \quad\Longleftrightarrow\quad
  \Obs(F)=0
  \quad\Longleftrightarrow\quad
  I_R=I_\Delta.
\]
\end{proposition}

\begin{proof}
The kernel of diagonal evaluation \(S\to A\) is \(I_\Delta\).  Passing
through \(S/I_R\) therefore identifies the displayed kernel with
\(I_\Delta/I_R\).  Since \(I_R\subseteq I_\Delta\),
\[
  I_\Delta/I_R=0
  \quad\Longleftrightarrow\quad
  I_R=I_\Delta.
\]
\end{proof}

\subsection{The fiber-product interpretation}
\label{sec:fiber-product}

Let
\[
  A=\C[x_1,\dots,x_n],\qquad
  B=\C[F_1,\dots,F_n]\subseteq A.
\]
Let \(X=\Spec A\), let \(Y=\Spec B\), and let
\[
  f_F\colon X\longrightarrow Y
\]
be the morphism induced by \(B\hookrightarrow A\).
The homomorphisms
\(\iota_x,\iota_y\colon A\rightrightarrows C_F\), defined by
\(\iota_x(a)=[a(x)]\) and \(\iota_y(a)=[a(y)]\), agree on \(B\) because
\([F_i(x)]=[F_i(y)]\) in \(C_F\).  Their common restriction equips
\(C_F\) with a \(B\)-algebra structure.

\begin{theorem}[Fiber-product interpretation]
\label{thm:collision-tensor-product}
There is a canonical \(B\)-algebra equivalence and its dual
affine-scheme equivalence
\[
  C_F\cong A\otimes_B A,
  \qquad
  R_F:=\Spec C_F\cong X\times_YX.
\]
Under the first equivalence, \(\bar\mu_F\) is tensor multiplication
\[
  A\otimes_BA\longrightarrow A,\qquad a\otimes a'\longmapsto aa'.
\]
Its dual morphism is the diagonal
\[
  \Delta_{X/Y}\colon X\longrightarrow X\times_YX.
\]
\end{theorem}

\begin{proof}
A map \(C_F\to T\) to a commutative \(\C\)-algebra \(T\) is equivalent
to a pair of maps \(f,g\colon A\rightrightarrows T\) satisfying
\[
  f(F_i)=g(F_i)\qquad(1\leq i\leq n).
\]
Since the \(F_i\) generate \(B\), this is precisely the universal
property of \(A\otimes_BA\).  Applying \(\Spec\) to this equivalence and
using
\[
  \Spec(A\otimes_BA)
  \cong
  \Spec A\times_{\Spec B}\Spec A
\]
gives the asserted fiber-product equivalence.
The statements about multiplication and the diagonal follow on pure
tensors.
\end{proof}

\[
  R_F(\C)
  =
  \{(u,v)\in X(\C)\times X(\C):F(u)=F(v)\}.
\]
Inside
\[
  X\times X=\Spec S,
\]
the ideals \(I_R\) and \(I_\Delta\) define, respectively,
\(X\times_YX\) and the image of \(\Delta_{X/Y}\).  The
contravariance of \(I\mapsto V(I)\) therefore gives the equivalences
\[
  I_R\subseteq I_\Delta
  \quad\Longleftrightarrow\quad
  \Delta_{X/Y}(X)\subseteq X\times_YX
\]
and
\[
  I_R=I_\Delta
  \quad\Longleftrightarrow\quad
  X\times_YX=\Delta_{X/Y}(X)
\]
as closed subschemes of \(X\times X\).  Thus the algebraic and
geometric gaps are two descriptions of the same scheme-theoretic
containment.

\subsection{The off-diagonal factor and secant projectors}
\label{sec:secant-projector}

For every polynomial map \(F\), define
\[
  I_{\mathrm{off}}:=I_R:I_\Delta,
  \qquad
  C_F^\circ:=S/I_{\mathrm{off}},
  \qquad
  R_F^\circ:=\Spec C_F^\circ.
\]
This definition does not require the diagonal to be clopen.

\begin{corollary}[Off-diagonal emptiness criterion]
\label{cor:off-diagonal-empty-criterion}
For every polynomial map \(F\), the following are equivalent:
\[
  R_F^\circ=\varnothing
  \quad\Longleftrightarrow\quad
  C_F^\circ=0
  \quad\Longleftrightarrow\quad
  \Obs(F)=0
  \quad\Longleftrightarrow\quad
  I_R=I_\Delta.
\]
\end{corollary}

\begin{proof}
\[
  R_F^\circ=\varnothing
  \quad\Longleftrightarrow\quad
  C_F^\circ=0
  \quad\Longleftrightarrow\quad
  I_R:I_\Delta=S.
\]
Moreover,
\[
  I_R:I_\Delta=S
  \quad\Longleftrightarrow\quad
  I_\Delta\subseteq I_R
  \quad\Longleftrightarrow\quad
  I_R=I_\Delta.
\]
The equivalence with \(\Obs(F)=0\) is
Proposition~\ref{prop:obstruction-kernel}.
\end{proof}

The idempotent mechanism behind a clopen diagonal is purely
commutative-algebraic.  Let \(C\) be a commutative ring, \(J\subseteq C\)
an ideal, \(\delta\in C\), and \(c\in C^\times\).  Suppose
\[
  \delta J=0,
  \qquad
  \delta-c\in J.
\]
Define
\[
  q:=1-c^{-1}\delta.
\]

\begin{theorem}[Abstract off-diagonal projector]
\label{thm:abstract-off-diagonal-projector}
The element \(q\) is idempotent and
\[
  J=Cq,
  \qquad
  \Ann_C(J)=C(1-q).
\]
In particular,
\[
  J=0\quad\Longleftrightarrow\quad q=0.
\]
\end{theorem}

\begin{proof}
The congruence \(\delta\equiv c\pmod J\) gives \(q\in J\).  For every
\(j\in J\),
\[
  qj=j-c^{-1}\delta j=j.
\]
Taking \(j=q\) proves \(q^2=q\); the same identity gives \(J=Cq\).
Finally,
\[
  aq=0
  \quad\Longleftrightarrow\quad
  a=a(1-q),
\]
so \(\Ann_C(J)=\Ann_C(Cq)=C(1-q)\).
\end{proof}

Whenever the theorem is applied with
\[
  C=C_F,
  \qquad
  J=\Obs(F),
\]
the identities \(J=C_Fq\) and \(C_F/J\cong A\) give
\[
  C_F
  \cong
  A
  \times
  C_F/C_F(1-q).
\]
Thus the first factor is the diagonal and the second is its
off-diagonal complement.

\subsection{Generic fibers and marked-root decomposition}
\label{sec:generic-galois-bridge}

Let \(F\) be dominant and generically finite, and let
\[
  A=\C[x_1,\dots,x_n],
  \qquad
  B=\C[F_1,\dots,F_n]\subseteq A.
\]
Let
\[
  K=\Frac(B),
  \qquad
  L=\Frac(A).
\]
Let
\[
  \theta\colon K\otimes_BA\longrightarrow L,
  \qquad
  \lambda\otimes a\longmapsto\lambda a.
\]

\begin{theorem}[Generic collision bridge]
\label{thm:generic-collision-base-change}
The map \(\theta\) is an isomorphism.  Consequently, there is a
\(K\)-algebra isomorphism
\[
  K\otimes_BC_F
  \cong
  L\otimes_KL,
\]
and the following square commutes:
\[
  \begin{array}{ccc}
    K\otimes_BC_F
      &\xrightarrow{\ 1\otimes\bar\mu_F\ }&
    K\otimes_BA\\
    {\scriptstyle\sim}\downarrow
      &&\downarrow{\scriptstyle\theta\ \sim}\\
    L\otimes_KL
      &\xrightarrow{\ \mu_L\ }&
    L.
  \end{array}
\]
\end{theorem}

\begin{proof}
Let \(U=B\setminus\{0\}\).  Since \(K=U^{-1}B\), localization gives
\[
  K\otimes_BA
  \xrightarrow{\sim}
  U^{-1}A,
  \qquad
  \frac{b}{u}\otimes a\longmapsto\frac{ba}{u}.
\]
Since \(A=\C[x_1,\dots,x_n]\) is a domain and \(B\subseteq A\),
\[
  B\setminus\{0\}\subseteq A\setminus\{0\}.
\]
Hence
\[
  U^{-1}A\hookrightarrow\Frac(A)=L,
  \qquad
  \frac{a}{u}\longmapsto\frac{a}{u}.
\]
Composing this injection with the displayed isomorphism gives
\(\theta\), so \(\theta\) is injective.

Since \(F\) is generically finite, \(L/K\) is a finite extension.  The
image of \(U^{-1}A\) in \(L\) is
\[
  K[x_1,\dots,x_n].
\]
Since each \(x_i\) is algebraic over \(K\), this \(K\)-domain is
finite-dimensional over \(K\), hence is a field.  Since it contains
\(A\), it contains \(\Frac(A)=L\); the reverse inclusion is already
given by its embedding into \(L\).  Therefore
\[
  U^{-1}A=L,
\]
and \(\theta\) is an isomorphism.

By Theorem~\ref{thm:collision-tensor-product}, identify
\(C_F\) with \(A\otimes_BA\).  There is a \(K\)-algebra isomorphism
\[
  \beta\colon
  K\otimes_B(A\otimes_BA)
  \xrightarrow{\sim}
  (K\otimes_BA)\otimes_K(K\otimes_BA)
\]
defined on pure tensors by
\[
  \beta\bigl(\lambda\otimes(a\otimes a')\bigr)
  =
  (\lambda\otimes a)\otimes(1\otimes a').
\]
Its inverse is
\[
  (\lambda\otimes a)\otimes(\lambda'\otimes a')
  \longmapsto
  \lambda\lambda'\otimes(a\otimes a').
\]
The balancing relations over \(B\) and \(K\) show that these formulas
are well-defined, and the two displayed maps are mutually inverse.
Composing \(\beta\) with \(\theta\otimes_K\theta\) gives
\[
  \Phi_F\colon K\otimes_BC_F\xrightarrow{\sim}L\otimes_KL.
\]

It remains to identify the diagonal map.  Under
\(C_F\cong A\otimes_BA\), the homomorphism \(\bar\mu_F\) is
multiplication.  Therefore, on a pure tensor,
\[
  \begin{aligned}
  (\mu_L\circ\Phi_F)
    \bigl(\lambda\otimes(a\otimes a')\bigr)
    &=(\lambda a)a'
     =\lambda aa',\\
  \bigl(\theta\circ(1\otimes\bar\mu_F)\bigr)
    \bigl(\lambda\otimes(a\otimes a')\bigr)
    &=\theta(\lambda\otimes aa')
     =\lambda aa'.
  \end{aligned}
\]
Pure tensors generate \(K\otimes_BC_F\), so the two homomorphisms are
equal.  This proves the commutative square.
\end{proof}

Suppose now that \(L/K\) is separable and has a power basis
\(L=K(\alpha)\).  Let \(m_\alpha(T)\in K[T]\) be the minimal polynomial
of \(\alpha\), and let \(h_\alpha(T)\in L[T]\) be determined by
\[
  m_\alpha(T)=(T-\alpha)h_\alpha(T)
  \qquad\text{in }L[T].
\]

\begin{theorem}[Marked-root tensor decomposition]
\label{thm:marked-root-tensor-decomposition}
Give \(L\otimes_KL\) its \(L\)-algebra structure through the left tensor
factor.  There is an \(L\)-algebra equivalence
\[
  L\otimes_KL
  \cong
  L\times L[T]/(h_\alpha),
\]
under which tensor multiplication \(L\otimes_KL\to L\) is projection
onto the first factor.
\end{theorem}

\begin{proof}
The power-basis presentation
\[
  L\cong K[T]/(m_\alpha)
\]
and base change from \(K\) to \(L\) give an \(L\)-algebra isomorphism
\[
  L\otimes_KL
  \xrightarrow{\sim}
  L[T]/(m_\alpha),
  \qquad
  1\otimes\alpha\longmapsto T.
\]
Since \(m_\alpha\) is separable,
\(m_\alpha'(\alpha)\neq0\).  Differentiating
\[
  m_\alpha(T)=(T-\alpha)h_\alpha(T)
\]
and evaluating at \(T=\alpha\) gives
\[
  h_\alpha(\alpha)=m_\alpha'(\alpha)\neq0.
\]
Thus \(T-\alpha\) does not divide \(h_\alpha\), so these two polynomials
are coprime in \(L[T]\).  The Chinese remainder theorem now gives
\[
  L[T]/(m_\alpha)
  \xrightarrow{\sim}
  L[T]/(T-\alpha)\times L[T]/(h_\alpha)
  \xrightarrow{\sim}
  L\times L[T]/(h_\alpha).
\]
Under the first isomorphism above, a pure tensor
\(\ell\otimes p(\alpha)\), with \(p\in K[T]\), corresponds to the class
of \(\ell p(T)\).  Tensor multiplication sends it to
\(\ell p(\alpha)\), which is evaluation at \(T=\alpha\).  Under the
Chinese remainder isomorphism, this evaluation map is projection onto
the factor \(L[T]/(T-\alpha)\cong L\).  Hence tensor multiplication is
the first projection.
\end{proof}

Together with Theorem~\ref{thm:generic-collision-base-change}, this gives
\[
  K\otimes_BC_F
  \cong
  L\times L[T]/(h_\alpha),
\]
where the first projection is the generic diagonal.

\subsection{Normalized covers and inertia}
\label{sec:normalized-covers-inertia}

Let \(L/K\) be finite and separable, and let \(N/K\) be its normal
closure, with a fixed embedding \(L\hookrightarrow N\).  Let
\[
  G=\Gal(N/K),
  \qquad
  H=\Gal(N/L).
\]

\begin{proposition}[Conjugate-embedding action]
\label{prop:conjugate-embedding-action}
Restriction induces a bijection
\[
  G/H\xrightarrow{\sim}\Hom_K(L,N),
  \qquad
  gH\longmapsto g|_L.
\]
The induced action of \(G\) on \(G/H\) is faithful; equivalently,
\[
  \core_G(H)=\bigcap_{g\in G}gHg^{-1}=1.
\]
\end{proposition}

\begin{proof}
For \(g_1,g_2\in G\),
\[
  g_1|_L=g_2|_L
  \quad\Longleftrightarrow\quad
  g_2^{-1}g_1\in H
  \quad\Longleftrightarrow\quad
  g_1H=g_2H.
\]
Every \(K\)-embedding \(L\hookrightarrow N\) extends to an element of
\(G\), so restriction gives the stated bijection.  Its action kernel is
\[
  \ker\!\left(G\curvearrowright G/H\right)
  =
  \bigcap_{g\in G}gHg^{-1}
  =
  \core_G(H).
\]
If \(\sigma\in\core_G(H)\), then
\[
  \sigma|_{g(L)}=\id_{g(L)}
  \quad(g\in G).
\]
Since \(N\) is the normal closure,
\[
  N=K\bigl(g(L):g\in G\bigr),
\]
and therefore \(\sigma=\id_N\).
\end{proof}

Let
\[
  Y=\Spec B,
  \qquad
  \overline X=\Norm_L(Y),
  \qquad
  Z=\Norm_N(Y).
\]
The marked embedding \(L\hookrightarrow N\) induces a canonical
commutative triangle
\[
  Z\xrightarrow{\pi}\overline X\xrightarrow{\bar\nu}Y,
  \qquad
  \nu=\bar\nu\circ\pi\colon Z\longrightarrow Y.
\]
Suppose that \(\bar\nu\) and \(\nu\) are finite, and fix an open
immersion over \(Y\)
\[
  j\colon X=\Spec A\hookrightarrow\overline X
\]
with \(\bar\nu\circ j=f_F\).  Denote this diagram and its deleted boundary
by
\[
  \mathcal D=
  \left(
    Z\xrightarrow{\pi}\overline X\xrightarrow{\bar\nu}Y,\,
    j\colon X\hookrightarrow\overline X
  \right),
  \qquad
  \partial X=\overline X\setminus j(X).
\]

Fix \(E\in Z^{(1)}\).  Let \(w_E\) be its divisorial
valuation, let \(b=\nu(E)\), and let \(v_b=w_E|_K\).
Let \(D_E\leq G\) and \(I_E\trianglelefteq D_E\) be its decomposition and
inertia groups.  Let
\[
  \mathcal Q_E:=D_E\backslash G/H.
\]
For \(\omega=D_EgH\in\mathcal Q_E\), let
\[
  u_\omega=(g^{-1}w_E)|_L,
  \qquad
  \operatorname{cent}(E,\omega)
  =\operatorname{cent}_{\overline X}(u_\omega),
\]
and let
\[
  i_E(\omega)=[I_E:I_E\cap gHg^{-1}].
\]
\begin{proposition}[Conjugate divisors and double cosets]
\label{prop:conjugate-divisor-double-cosets}
The assignments above are well-defined, and restriction of valuations
induces a bijection
\[
  \mathcal Q_E
  \xrightarrow{\sim}
  \{\text{prime divisors of }\overline X\text{ above }b\},
  \qquad
  \omega\longmapsto\operatorname{cent}(E,\omega).
\]
\end{proposition}

\begin{proof}
The group \(G\) acts transitively on the extensions of \(v_b\) to \(N\),
and the stabilizer of \(w_E\) is \(D_E\).  Restricting such a valuation
to \(L=N^H\) identifies precisely those extensions whose indices differ
on the right by an element of \(H\).  The resulting orbits are therefore
indexed by \(D_E\backslash G/H\), giving the stated bijection.

It remains to check independence of representatives.  If
\(D_Eg'H=D_EgH\), then \(g'=dgh\) for some \(d\in D_E\) and \(h\in H\).
The element \(d\) fixes \(w_E\), while \(h\) fixes \(L\) pointwise, so
\[
  (g'^{-1}w_E)|_L=(g^{-1}w_E)|_L.
\]
Moreover, \(I_E\trianglelefteq D_E\), and conjugation by \(d\) identifies
\[
  I_E\cap g'Hg'^{-1}
  \quad\text{with}\quad
  I_E\cap gHg^{-1}.
\]
Thus \(u_\omega\), \(\operatorname{cent}(E,\omega)\), and
\(i_E(\omega)\) are
well-defined.
\end{proof}

We use the notation
\[
  e\bigl(\operatorname{cent}(E,\omega)/b\bigr):=e(u_\omega/v_b).
\]
Since the residue characteristic is zero,
\[
  \Ram^{(1)}(Z/Y)
  :=
  \{E\in Z^{(1)}:e(w_E/v_b)>1\}
  =
  \{E\in Z^{(1)}:I_E\neq1\}.
\]

\begin{proposition}[Detection of inertia on a conjugate sheet]
\label{prop:core-free-sheet-action}
If \(I_E\neq1\), then there is an \(\omega\in\mathcal Q_E\) with
\[
  i_E(\omega)>1.
\]
\end{proposition}

\begin{proof}
Suppose, to the contrary, that \(i_E(\omega)=1\) for every
\(\omega\in\mathcal Q_E\).  For each \(g\in G\),
\[
  1=i_E(D_EgH)=[I_E:I_E\cap gHg^{-1}],
\]
so
\[
  I_E=I_E\cap gHg^{-1}\subseteq gHg^{-1}
  \qquad\text{for every }g\in G.
\]
Consequently,
\[
  I_E
  \subseteq
  \bigcap_{g\in G}gHg^{-1}
  =
  \core_G(H)
  =
  1,
\]
contrary to \(I_E\neq1\).
\end{proof}

\begin{proposition}[Intermediate ramification index]
\label{prop:intermediate-ramification-index}
For \(\omega=D_EgH\in\mathcal Q_E\),
\[
  e\bigl(\operatorname{cent}(E,\omega)/b\bigr)
  =
  i_E(\omega)
  =
  [I_E:I_E\cap gHg^{-1}].
\]
\end{proposition}

\begin{proof}
Let \(w_E\) and \(v_b\) be the divisorial valuations associated with
\(E\) and \(b\), and let
\[
  w_g=g^{-1}w_E,
  \qquad
  u_\omega=w_g|_L.
\]
The valuation \(u_\omega\) has center \(\operatorname{cent}(E,\omega)\) on
\(\overline X\).  Since
the residue fields have characteristic zero,
\[
  e(w_g/v_b)=|I_E|.
\]
The inertia group of \(w_g\) over \(u_\omega\) is
\[
  g^{-1}I_Eg\cap H,
\]
and hence
\[
  e(w_g/u_\omega)
  =
  |g^{-1}I_Eg\cap H|
  =
  |I_E\cap gHg^{-1}|.
\]
Multiplicativity of ramification indices in the field tower
\[
  K\subseteq L\subseteq N
\]
therefore gives
\[
  e(u_\omega/v_b)
  =
  \frac{e(w_g/v_b)}{e(w_g/u_\omega)}
  =
  [I_E:I_E\cap gHg^{-1}],
\]
which is the asserted formula.
\end{proof}

\begin{proposition}[Visibility-to-boundary]
\label{prop:visibility-to-boundary}
Suppose that \(X\to Y\) is \'etale.  For \(E\in Z^{(1)}\) and
\(\omega\in\mathcal Q_E\),
\[
  i_E(\omega)>1
  \quad\Longrightarrow\quad
  \operatorname{cent}(E,\omega)\in\partial X.
\]
\end{proposition}

\begin{proof}
Suppose that \(\operatorname{cent}(E,\omega)\in j(X)\), and let
\(x\in X\) be the corresponding point under the open immersion \(j\).
The local homomorphism
\[
  \mathcal O_{Y,b}
  \longrightarrow
  \mathcal O_{X,x}
\]
is \'etale, so
\(e(\operatorname{cent}(E,\omega)/b)=1\).  Proposition
\ref{prop:intermediate-ramification-index} now gives
\[
  i_E(\omega)
  =
  e\bigl(\operatorname{cent}(E,\omega)/b\bigr)
  =
  1.
\]
Taking the contrapositive gives the result.
\end{proof}

\begin{definition}[Hidden-inertia locus]
\label{def:no-hidden-inertia}
\[
  \mathcal H(\mathcal D)
  :=
  \left\{
  E\in\Ram^{(1)}(Z/Y):
  \begin{array}{l}
    \exists\omega\in\mathcal Q_E,\ i_E(\omega)>1,\\
    \forall\omega\in\mathcal Q_E,\quad
      i_E(\omega)>1
      \Longrightarrow
      \operatorname{cent}(E,\omega)\in\partial X
  \end{array}
  \right\}.
\]
Thus \(E\in\mathcal H(\mathcal D)\) precisely when nontrivial inertia
is visible on a conjugate sheet but every visible center belongs to the
deleted boundary.
\end{definition}

\begin{proposition}[Divisorial ramification is hidden from the \'etale sheet]
\label{prop:etale-sheet-hidden-inertia}
Suppose that \(N/K\) is the normal closure of \(L/K\), the residue
characteristic is zero, and \(X\to Y\) is \'etale.  Then
\[
  \mathcal H(\mathcal D)=\Ram^{(1)}(Z/Y).
\]
Equivalently, every divisorial ramification orbit has a conjugate sheet on
which its inertia acts nontrivially, but every such center lies in the
deleted boundary.
\end{proposition}

\begin{proof}
The inclusion
\(\mathcal H(\mathcal D)\subseteq\Ram^{(1)}(Z/Y)\) is part of the
definition.  Conversely, let \(E\in\Ram^{(1)}(Z/Y)\).  In residue
characteristic zero, \(I_E\neq1\).  Since \(N\) is the normal closure of
\(L/K\), Proposition~\ref{prop:conjugate-embedding-action} gives
\(\core_G(H)=1\), and Proposition~\ref{prop:core-free-sheet-action}
supplies an \(\omega\in\mathcal Q_E\) with \(i_E(\omega)>1\).
For every \(\omega'\in\mathcal Q_E\), Proposition
\ref{prop:visibility-to-boundary} gives
\[
  i_E(\omega')>1
  \quad\Longrightarrow\quad
  \operatorname{cent}(E,\omega')\in\partial X.
\]
These are exactly the two conditions defining
\(E\in\mathcal H(\mathcal D)\).
\end{proof}

\subsection{Automorphism detection}
\label{sec:automorphism-detection}

\begin{theorem}[Automorphism criterion]
\label{thm:automorphism-criterion}
For an arbitrary polynomial self-map
\(F\colon\A^n_{\C}\to\A^n_{\C}\), the following are equivalent:
\begin{enumerate}[label=\textup{(\arabic*)}]
  \item \(F\) is a polynomial automorphism;
  \item \(F\) is injective on \(\A^n(\C)\);
  \item \(I_R(F)=I_\Delta\);
  \item \(\Obs(F)=0\).
\end{enumerate}
\end{theorem}

\begin{proof}
Suppose first that \(F\) has polynomial inverse
\(G=(G_1,\dots,G_n)\).  Then
\[
  x_i-y_i
  =
  G_i(F(x))-G_i(F(y))
  \in I_R
\]
for every \(i\), so \(I_\Delta\subseteq I_R\).  Proposition
\ref{prop:collision-contained-in-diagonal} gives equality.  Thus
\textup{(1)} implies \textup{(3)}.

Suppose \(I_R=I_\Delta\), and let
\(\operatorname{ev}_{(a,b)}\colon S\to\C\) be evaluation at
\((a,b)\).  If \(F(a)=F(b)\), then
\[
  I_R\subseteq\ker(\operatorname{ev}_{(a,b)}).
\]
Therefore \(x_i-y_i\in I_\Delta=I_R\) gives \(a_i-b_i=0\) for every
\(i\).  Thus \(a=b\), proving
\textup{(3)}\(\Rightarrow\)\textup{(2)}.  The
Ax--Grothendieck theorem \cite{Ax1969} gives
\textup{(2)}\(\Rightarrow\)\textup{(1)}.  Finally,
Proposition~\ref{prop:obstruction-kernel} gives
\textup{(3)}\(\Longleftrightarrow\)\textup{(4)}.
\end{proof}

\begin{corollary}[The Jacobian conjecture as obstruction vanishing]
\label{cor:jacobian-conjecture-as-vanishing}
For every \(n\geq1\), the \(n\)-dimensional Jacobian conjecture is
equivalent to each of the statements
\[
  \forall F\colon\A^n_{\C}\to\A^n_{\C},
  \qquad
  \det JF\in\C^\times
  \quad\Longrightarrow\quad
  \Obs(F)=0,
\]
and
\[
  \forall F\colon\A^n_{\C}\to\A^n_{\C},
  \qquad
  \det JF\in\C^\times
  \quad\Longrightarrow\quad
  I_R(F)=I_\Delta.
\]
\end{corollary}

\begin{proof}
Apply Theorem~\ref{thm:automorphism-criterion} to every Keller map.
Its equivalent conditions \textup{(1)}, \textup{(3)}, and \textup{(4)}
give the two displayed formulations.
\end{proof}

\subsection{Generic degree one}

\begin{corollary}[Generic degree-one descent]
\label{cor:generic-degree-one-descent}
If
\[
  L=K
  \qquad\text{and}\qquad
  A\text{ is flat over }B,
\]
then
\[
  \Obs(F)=0.
\]
\end{corollary}

\begin{proof}
When \(L=K\), Theorem~\ref{thm:generic-collision-base-change}
identifies the base-changed diagonal with tensor multiplication:
\[
  \ker(1\otimes\bar\mu_F)
  \cong
  \ker\!\left(
    K\otimes_KK\xrightarrow{\mu_K}K
  \right)
  =
  0.
\]
The fiber-product identification of
Theorem~\ref{thm:collision-tensor-product} is
\[
  C_F\xrightarrow{\sim}A\otimes_BA.
\]
Since \(A\) is flat over \(B\), the tensor product \(A\otimes_BA\), and
therefore \(C_F\), is flat over \(B\).  The localization map
\[
  \iota_C\colon C_F\hookrightarrow K\otimes_BC_F
\]
is therefore injective, and
\[
  \iota_C\bigl(\ker\bar\mu_F\bigr)
  \subseteq
  \ker(1\otimes\bar\mu_F)
  =
  0.
\]
Thus \(\ker\bar\mu_F=0\).  Proposition
\ref{prop:obstruction-kernel} identifies
\[
  \ker(\bar\mu_F)=\Obs(F)=I_\Delta/I_R,
\]
so \(\Obs(F)=0\).
\end{proof}

\begin{leanfiles}
\leanpath{CollisionIdeals/CollisionDiagonal.lean},
\leanpath{CollisionIdeals/OffDiagonalScheme.lean},
\leanpath{CollisionIdeals/Secant.lean},
\leanpath{CollisionIdeals/GenericCollisionDiagonal.lean},
\leanpath{CollisionIdeals/KellerCollisionModel.lean}, and
\leanpath{CollisionIdeals/JacobianConjecture.lean}.  The Lean form of
Theorem~\ref{thm:generic-collision-base-change} takes surjectivity of
\(\theta\) as a parameter; the proof above derives it from generic
finiteness.
\end{leanfiles}

% The planar specialization is retained in main.tex and reserved for Paper II.

\section{\texorpdfstring{The cubic \(S_3\) obstruction in dimension three}
  {The cubic S3 obstruction in dimension three}}
\label{sec:cubic-s3-obstruction}

We now specialize the dimension-independent constructions of
Section~\ref{sec:general-construction} to dimension three and isolate the
separable nonnormal cubic case.

Let \(S_3=\operatorname{Perm}(\{1,2,3\})\) be the symmetric group on
three letters.

\begin{proposition}[Dimension-three counterexample criterion]
\label{prop:dimension-three-counterexample-criterion}
For a polynomial map
\[
  F\colon\A^3_{\C}\longrightarrow\A^3_{\C},
\]
consider the following statements:
\begin{enumerate}[label=\textup{(\arabic*)}]
  \item \(\det JF\in\C^\times\);
  \item \(F\) is not a polynomial automorphism;
  \item \(R_F^\circ\neq\varnothing\).
\end{enumerate}
Statements \textup{(2)} and \textup{(3)} are equivalent.  Consequently,
\(F\) is a dimension-three Jacobian-conjecture counterexample exactly
when \textup{(1)} and \textup{(3)} hold.
\end{proposition}

\begin{proof}
By Corollary~\ref{cor:off-diagonal-empty-criterion} and
Theorem~\ref{thm:automorphism-criterion},
\[
  R_F^\circ=\varnothing
  \quad\Longleftrightarrow\quad
  I_R=I_\Delta
  \quad\Longleftrightarrow\quad
  F\in\Aut_{\mathrm{poly}}(\A^3_{\C}).
\]
Negating this equivalence proves
\textup{(2)}\(\Longleftrightarrow\)\textup{(3)}.  Combining this with
\textup{(1)} gives the final assertion.
\end{proof}

\subsection{The cubic extension and its normal closure}
\label{sec:cubic-extension}

Let
\[
  F=(F_1,F_2,F_3)\colon\A^3_{\C}\longrightarrow\A^3_{\C}
\]
be dominant and generically finite.  Let
\[
  A=\C[x_1,x_2,x_3],
  \qquad
  B=\C[F_1,F_2,F_3]\subseteq A,
  \qquad
  K=\Frac(B),
  \qquad
  L=\Frac(A),
\]
and let
\[
  \theta\colon K\otimes_BA\longrightarrow L,
  \qquad
  \lambda\otimes a\longmapsto\lambda a,
\]
be the isomorphism of
Theorem~\ref{thm:generic-collision-base-change}.  Suppose that \(L/K\)
is separable and nonnormal, and
choose a power basis with generator \(\alpha\) and degree three:
\[
  L=K(\alpha),
  \qquad
  [L:K]=3.
\]
If \(m_\alpha\in K[T]\) is the minimal polynomial of \(\alpha\), let
\(h_\alpha\in L[T]\) be determined by
\[
  m_\alpha(T)=(T-\alpha)h_\alpha(T)
  \qquad\text{in }L[T].
\]
Let \(N/K\) be the normal closure of \(L/K\), containing \(L\) through
the marked inclusion.  Let
\[
  G=\Gal(N/K),
  \qquad
  H=\Gal(N/L).
\]
Recall from Section~\ref{sec:normalized-covers-inertia} that
\[
  G/H\xrightarrow{\sim}\Hom_K(L,N),
  \qquad
  gH\longmapsto g|_L,
  \qquad
  \core_G(H)=1.
\]
The final equality follows because \(N\) is the normal closure of
\(L/K\); equivalently, the action of \(G\) on the conjugate embeddings
of \(L\) is faithful.

For \(Y=\Spec B\), the finite-normalization construction of
Section~\ref{sec:normalized-covers-inertia} gives
\[
  Z=\Norm_N(Y)\longrightarrow
  \overline X=\Norm_L(Y)\longrightarrow Y.
\]
If \(E\in Z^{(1)}\), with decomposition group \(D_E\) and inertia group
\(I_E\), then its conjugate centers are indexed by
\(\mathcal Q_E=D_E\backslash G/H\).  For
\(\omega=D_EgH\), Proposition~\ref{prop:intermediate-ramification-index}
states
\[
  e\bigl(\operatorname{cent}(E,\omega)/\nu(E)\bigr)
  =
  i_E(\omega)
  =
  [I_E:I_E\cap gHg^{-1}].
\]

The following proposition is the standard residual-factor description
of a separable nonnormal cubic extension.  We include its short proof
to fix the marked identification with \(N\).

\begin{proposition}[Cubic residual-factor identification]
\label{prop:cubic-tensor-decomposition}
The polynomial \(h_\alpha\) is irreducible over \(L\), and there is an
\(L\)-algebra isomorphism
\[
  \varphi\colon L[T]/(h_\alpha)\xrightarrow{\sim}_L N.
\]
Consequently,
\[
  \Psi\colon L\otimes_KL\xrightarrow{\sim}_L L\times N,
  \qquad
  \operatorname{pr}_1\circ\Psi=\mu_L.
\]
\end{proposition}

\begin{proof}
Write
\[
  m_\alpha(T)=T^3+a_2T^2+a_1T+a_0.
\]
If the quadratic \(h_\alpha\) were reducible over \(L\), it would have a
root \(\beta\in L\).  Separability gives \(\beta\neq\alpha\), and the
third root would be
\[
  \gamma=-a_2-\alpha-\beta\in L.
\]
Thus \(m_\alpha\) would split over \(L\).  Since
\(L=K(\alpha)\), the extension \(L/K\) would then be the splitting
field of the separable polynomial \(m_\alpha\), hence normal.  This
contradicts the assumed nonnormality, so \(h_\alpha\) is irreducible.

Let \(\beta\) denote the image of \(T\) in \(L[T]/(h_\alpha)\).  This
field contains the three roots
\[
  \alpha,\qquad \beta,\qquad -a_2-\alpha-\beta
\]
of \(m_\alpha\), and is generated over \(L\) by \(\beta\).  It is
therefore the splitting field of \(m_\alpha\), which is the normal
closure \(N\) of \(L/K\).  Choose the resulting \(L\)-algebra
isomorphism
\[
  \varphi\colon L[T]/(h_\alpha)\xrightarrow{\sim}_L N.
\]

Let
\[
  \Phi_\alpha\colon
  L\otimes_KL
  \xrightarrow{\sim}_L
  L\times L[T]/(h_\alpha)
\]
be the marked-root decomposition of
Theorem~\ref{thm:marked-root-tensor-decomposition}.  It satisfies
\[
  \operatorname{pr}_1\circ\Phi_\alpha=\mu_L.
\]
Composing its second factor with \(\varphi\), let
\[
  \Psi=(\id_L\times\varphi)\circ\Phi_\alpha.
\]
Since \(\id_L\times\varphi\) leaves the first factor unchanged,
\(\operatorname{pr}_1\circ\Psi=\mu_L\).
\end{proof}

\begin{theorem}[Galois group of a nonnormal cubic extension]
\label{thm:index-three-core-free-s3}
\[
  \Gal(N/K)\cong S_3.
\]
\end{theorem}

\begin{proof}
Proposition~\ref{prop:cubic-tensor-decomposition} identifies \(N\) with
the degree-two extension \(L[T]/(h_\alpha)\) of \(L\).  Hence
\[
  |H|=[N:L]=2.
\]
The normal-closure identities recalled above give
\[
  \core_G(H)=1,
  \qquad
  [G:H]=[L:K]=3.
\]
Hence the action of \(G\) on the three cosets of \(H\) gives
\[
  \rho\colon G\longrightarrow\operatorname{Perm}(G/H)\cong S_3,
  \qquad
  \ker\rho=\core_G(H)=1.
\]
Core-freeness therefore makes \(\rho\) injective, so
\[
  |G|\leq |S_3|=6.
\]
The index formula gives
\[
  |G|=[G:H]|H|=3\cdot2=6.
\]
Thus \(|G|=6\), and the injective homomorphism \(\rho\) is an
isomorphism \(G\cong S_3\).
\end{proof}

\subsection{The cubic off-diagonal factor}

Retain the cubic data of Section~\ref{sec:cubic-extension}, and suppose
in addition that \(F\) is Keller.  Thus
\[
  B=\C[F_1,F_2,F_3]\subseteq A=\C[x_1,x_2,x_3],
  \qquad
  K=\Frac(B)\subseteq L=\Frac(A)=K(\alpha)\subseteq N,
\]
where \(L/K\) is separable, nonnormal, and of degree three, and \(N/K\)
is its normal closure.  Let
\[
  \theta\colon K\otimes_BA\xrightarrow{\sim}L,
  \qquad
  \lambda\otimes a\longmapsto\lambda a.
\]

\begin{corollary}[Generic cubic \(S_3\) collision]
\label{cor:cubic-s3-collision}
There is a \(K\)-algebra isomorphism
\[
  \Psi_F\colon
  K\otimes_BC_F\xrightarrow{\sim}_K L\times N,
  \qquad
  \operatorname{pr}_1\circ\Psi_F
  =
  \theta\circ(1\otimes\bar\mu_F).
\]
It satisfies
\[
  \Psi_F\!\left(
    \ker\bigl(\theta\circ(1\otimes\bar\mu_F)\bigr)
  \right)
  =
  0\times N,
  \qquad
  \Gal(N/K)\cong S_3.
\]
Consequently,
\[
  \Obs(F)\neq0,
  \qquad
  I_R\subsetneq I_\Delta,
  \qquad
  R_F^\circ\neq\varnothing,
  \qquad
  F\notin\Aut_{\mathrm{poly}}(\A^3_{\C}).
\]
\end{corollary}

\begin{proof}
Theorem~\ref{thm:generic-collision-base-change} gives a \(K\)-algebra
isomorphism
\[
  \Phi_F\colon K\otimes_BC_F\xrightarrow{\sim}_K L\otimes_KL
\]
such that
\[
  \mu_L\circ\Phi_F
  =
  \theta\circ(1\otimes\bar\mu_F).
\]
Proposition~\ref{prop:cubic-tensor-decomposition} gives
\[
  \Psi\colon L\otimes_KL\xrightarrow{\sim}_L L\times N,
  \qquad
  \operatorname{pr}_1\circ\Psi=\mu_L.
\]
After restricting \(\Psi\) from \(L\) to \(K\), let
\[
  \Psi_F=\Psi\circ\Phi_F.
\]
It follows that
\[
  \operatorname{pr}_1\circ\Psi_F
  =
  \theta\circ(1\otimes\bar\mu_F),
  \qquad
  \Psi_F\!\left(
    \ker\bigl(\theta\circ(1\otimes\bar\mu_F)\bigr)
  \right)
  =
  0\times N\neq0.
\]

Suppose that the diagonal-evaluation homomorphism
\(\bar\mu_F\colon C_F\to A\) were injective.  Since \(K\) is a
localization of \(B\), it is flat over
\(B\), so \(1\otimes\bar\mu_F\) would also be injective.
Theorem~\ref{thm:generic-collision-base-change} shows that \(\theta\)
is an isomorphism.  Therefore
\(\theta\circ(1\otimes\bar\mu_F)\) would be injective, contradicting
the nonzero kernel \(0\times N\) above.  Hence
\(\ker(\bar\mu_F)\neq0\), and
Proposition~\ref{prop:obstruction-kernel} gives
\(\Obs(F)\neq0\).  Proposition
\ref{prop:collision-contained-in-diagonal} shows that
\(I_R\subseteq I_\Delta\) strict, and
Corollary~\ref{cor:off-diagonal-empty-criterion} gives
\(R_F^\circ\neq\varnothing\).

Finally, Theorem~\ref{thm:index-three-core-free-s3} identifies
\(\Gal(N/K)\) with \(S_3\).  If \(F\) were a polynomial automorphism,
Theorem~\ref{thm:automorphism-criterion} would give
\(\Obs(F)=0\), contrary to \(\Obs(F)\neq0\).  Since \(F\) is Keller, it is
a dimension-three Jacobian-conjecture counterexample.
\end{proof}

For the same cubic map and field tower, let
\[
  X=\Spec A,
  \qquad
  Y=\Spec B,
  \qquad
  G=\Gal(N/K),
  \qquad
  H=\Gal(N/L).
\]
Suppose there is a diagram over \(Y\)
\[
  \mathcal D=
  \left(
    Z=\Norm_N(Y)\xrightarrow{\pi}
    \overline X=\Norm_L(Y)\xrightarrow{\bar\nu}Y,\,
    j\colon X\hookrightarrow\overline X
  \right)
\]
in which \(\pi\) and \(\bar\nu\) are finite, \(j\) is an open immersion,
\(\bar\nu\circ j=f_F\), and \(f_F\colon X\to Y\) is \'etale.  Let
\[
  \partial X=\overline X\setminus j(X),
  \qquad
  G\cong S_3
\]
by Theorem~\ref{thm:index-three-core-free-s3}.

\begin{corollary}[Hidden inertia in the \(S_3\)-normalization]
\label{cor:cubic-hidden-inertia}
\[
  \mathcal H(\mathcal D)=\Ram^{(1)}(Z/Y).
\]
\end{corollary}

\begin{proof}
Apply Proposition~\ref{prop:etale-sheet-hidden-inertia} to the displayed
normalization diagram.
\end{proof}

\begin{remark}[Ordered roots]
Let
\[
  m_\alpha(T)=T^3+a_2T^2+a_1T+a_0
\]
have roots \(\alpha_1,\alpha_2,\alpha_3\) in \(N\).  For distinct
\(i,j,k\),
\[
  \alpha_k=-a_2-\alpha_i-\alpha_j,
  \qquad
  K(\alpha_i,\alpha_j)
  =
  K(\alpha_1,\alpha_2,\alpha_3)
  =
  N.
\]
Thus the off-diagonal factor, which marks two distinct roots, is already
the ordered-root normal closure.
\end{remark}

\begin{leanfiles}
\leanpath{CollisionIdeals/ComplexThree/CubicGaloisGroup.lean},
\leanpath{CollisionIdeals/ComplexThree/S3Collision.lean},
\leanpath{CollisionIdeals/ComplexThree/OffDiagonal.lean},
\leanpath{CollisionIdeals/ComplexThree/CubicS3.lean}, and
\leanpath{CollisionIdeals/ComplexThree/JacobianConjecture.lean}.  The
Lean cubic witness takes the residual-field isomorphism and
\(H\neq1\) as explicit inputs; Proposition~\ref{prop:cubic-tensor-decomposition}
derives them here from nonnormality.  Hidden inertia is formalized
pointwise at a supplied ramified divisor.
\end{leanfiles}

\begin{remark}[Exact relation with classical complex \(JC(3)\)]
\label{rem:exact-relation-jc3}
The traditional complex three-dimensional Jacobian conjecture asserts that
every polynomial map
\[
  F\colon\A^3_{\C}\longrightarrow\A^3_{\C},
  \qquad
  \det JF\in\C^\times,
\]
is a polynomial automorphism; equivalently, every such map satisfies
\(\Obs(F)=0\).  Corollary~\ref{cor:cubic-s3-collision} proves the following
implication under a prescribed function-field hypothesis:
\[
  \begin{gathered}
    \det JF\in\C^\times,\qquad [L:K]=3,\qquad
    L/K\text{ nonnormal}\\
    \Longrightarrow\quad \Obs(F)\neq0
    \quad\Longrightarrow\quad
    F\text{ is a counterexample to }JC(3).
  \end{gathered}
\]
Consequently, the existence of one such Keller map would disprove
\(JC(3)\), while \(JC(3)\) implies that no such map exists.  This section
does not construct such a map, and it does not assert that every possible
counterexample to \(JC(3)\) has generic degree three.  It identifies the
generic collision geometry of this cubic nonnormal subclass.
\end{remark}

\section{Conclusion}
\label{sec:conclusion}

The obstruction ideal \(\Obs(F)\) vanishes exactly when the
scheme-theoretic collision relation \(X\times_YX\) equals the diagonal.
In dimension two there is one endgame with several exact formulations.
PlanarBoundaryCoherence expresses it through disappearance of the whole
deleted boundary.  PlanarRamificationRigidity, moving-sheet coverage, and
boundary separation express its height-one premise in differential,
valuation-theoretic, and finite-group language.  Only after that premise is
established do purity
and finite-\'etale rigidity make the normal closure trivial and force
\(\Obs(F)=0\).  The image of the deleted boundary is exactly the
nonproperness set, and the boundary divisor-class argument proves the
normal/Galois case unconditionally.  Universal validity of the equivalent
planar criteria is exactly \(JC(2)\), not a result asserted here.  In
the cubic dimension-three class, by contrast, the nonzero generic factor is
the ordered-root \(S_3\)-normal closure.
Both statements arise from the same collision, off-diagonal, and
conjugate-sheet constructions; their point is to locate the obstruction,
respectively, at the normalization boundary and on the generic
off-diagonal sheet.

Under PlanarBoundaryCoherence or PlanarRamificationRigidity in dimension
two, under the conditions of Section~\ref{sec:cubic-s3-obstruction} in
dimension three, and with
\(\Ram^{(1)}(Z/Y)\neq\varnothing\) in the cubic normalization, the
contrast is
\[
  \begin{array}{c@{\qquad}|@{\qquad}c}
    \hline
    \text{\rm planar rigidity criteria}
      & \text{\rm nonnormal cubic class in dimension three}\\ \hline
    \operatorname{PBC}(F)\text{\rm\ or }\operatorname{PRR}(F)
      & \text{\rm hidden inertia in the \(S_3\)-cover realized}\\
    \Obs(F)=0 & \Obs(F)\neq0\\
    I_R=I_\Delta & I_R\subsetneq I_\Delta\\
    F\text{\rm\ is a polynomial automorphism}
      & F\text{\rm\ is not a polynomial automorphism}.\\
    \hline
  \end{array}
\]

\paragraph{Motivating example.}
The explicit three-dimensional polynomial counterexample announced by
Levent Alpöge \cite{Alpoge2026} and independently formalized in
Isabelle/HOL by Ramos, Hulak, and de Queiroz
\cite{AFPJacobianCounterexample} motivated the cubic \(S_3\) comparison.
The formal verification establishes a
constant nonzero Jacobian determinant and distinct rational points with
the same image.  That particular map is not used in the arguments above:
the dimension-three result is the structural implication from the
stated cubic and marked-normal-closure conditions.

The formalization uses Lean~4 and Mathlib \cite{Lean4,Mathlib}.

\paragraph{Acknowledgment.}
OpenAI Codex (GPT-5.6 Sol) assisted with Lean proof engineering,
mathematical and bibliographic checking, and manuscript editing.  The
author reviewed and takes responsibility for all statements, proofs,
citations, and formalization code.

\bibliographystyle{alpha}
\bibliography{references}

\end{document}
